\section{Численное дифференцирование}
\label{sec:numerical-differentiation}

\subsection{Постановка задачи}

Численное дифференцирование применяется, когда функция задана таблицей значений
или её аналитическая производная неудобна для вычисления. Производная заменяется
линейной комбинацией значений функции в соседних узлах:
\[
    f^{(m)}(x_0)\approx\sum_{j}a_jf(x_j).
\]
Коэффициенты выбираются так, чтобы формула была точна на полиномах максимально
возможной степени. Обычно они выводятся из разложения Тейлора или
дифференцирования интерполяционного многочлена.

\subsection{Формулы для первой производной}

На равномерной сетке с шагом $h$:
\[
    f'(x)=\frac{f(x+h)-f(x)}{h}+O(h)
\]
--- правая разность;
\[
    f'(x)=\frac{f(x)-f(x-h)}{h}+O(h)
\]
--- левая разность;
\[
    f'(x)=\frac{f(x+h)-f(x-h)}{2h}+O(h^2)
\]
--- центральная разность.

Формула второго порядка на границе:
\[
    f'(x)=\frac{-3f(x)+4f(x+h)-f(x+2h)}{2h}+O(h^2).
\]
Формула четвёртого порядка во внутренней точке:
\[
    f'(x)=\frac{f(x-2h)-8f(x-h)+8f(x+h)-f(x+2h)}{12h}
    +O(h^4).
\]

\subsection{Вторая и высшие производные}

Стандартная центральная формула для второй производной:
\[
    f''(x)=\frac{f(x-h)-2f(x)+f(x+h)}{h^2}+O(h^2).
\]
Формула четвёртого порядка:
\[
    f''(x)=\frac{-f(x+2h)+16f(x+h)-30f(x)+16f(x-h)-f(x-2h)}{12h^2}
    +O(h^4).
\]

Для смешанной производной функции двух переменных:
\[
\frac{\partial^2 f}{\partial x\partial y}(x,y)
\approx
\frac{f(x+h,y+k)-f(x+h,y-k)-f(x-h,y+k)+f(x-h,y-k)}{4hk}.
\]

\subsection{Вывод методом неопределённых коэффициентов}

Пусть требуется формула
\[
    f'(x_0)\approx\frac1h\sum_{j=-r}^{s}a_jf(x_0+jh).
\]
Подставим ряды Тейлора:
\[
    f(x_0+jh)=
    \sum_{k=0}^{\infty}\frac{(jh)^k}{k!}f^{(k)}(x_0).
\]
Для точности порядка $p$ коэффициенты должны удовлетворять моментным условиям
\[
    \sum_j a_j=0,
    \qquad
    \sum_j ja_j=1,
\]
\[
    \sum_j j^ka_j=0,
    \qquad k=2,\ldots,p.
\]
Решение этой линейной системы даёт разностную формулу требуемого порядка.

\subsection{Неравномерная сетка}

При неравных шагах используют интерполяционный многочлен Лагранжа или Ньютона и
дифференцируют его. Для трёх точек $x_{i-1}<x_i<x_{i+1}$ коэффициенты зависят от
обоих шагов
\[
    h_-=x_i-x_{i-1},
    \qquad
    h_+=x_{i+1}-x_i.
\]
В частности,
\[
\begin{aligned}
f'(x_i)\approx
&-\frac{h_+}{h_-(h_-+h_+)}f_{i-1}
+\frac{h_+-h_-}{h_-h_+}f_i\\
&+\frac{h_-}{h_+(h_-+h_+)}f_{i+1}.
\end{aligned}
\]
Формула имеет второй порядок при регулярном сгущении сетки.

\subsection{Погрешность и выбор шага}

Полная ошибка состоит из погрешности усечения и округления. Для центральной
разности первой производной:
\[
    E(h)\approx C_1h^2+C_2\frac{\varepsilon_{\mathrm{mach}}}{h}.
\]
При слишком большом $h$ доминирует ошибка аппроксимации, при слишком малом ---
катастрофическая потеря значащих цифр при вычитании близких чисел. Оптимальный
шаг по порядку величины:
\[
    h_{\mathrm{opt}}\sim\varepsilon_{\mathrm{mach}}^{1/3}
    \quad\text{для центральной разности первой производной}.
\]
Для второй производной баланс ещё жёстче, поскольку ошибка округления делится на
$h^2$.

Численное дифференцирование является некорректной по Адамару операцией: малый
высокочастотный шум в функции может давать большую ошибку производной.

\subsection{Экстраполяция Ричардсона}

Пусть приближение имеет разложение
\[
    D(h)=f'(x)+Ch^p+O(h^{p+1}).
\]
Тогда комбинация
\[
    D_{\mathrm{ext}}=
    \frac{2^pD(h/2)-D(h)}{2^p-1}
\]
устраняет главный член ошибки и повышает порядок точности. Метод можно применять
последовательно, строя таблицу экстраполяции.

\subsection{Дифференцирование зашумлённых данных}

Прямые разности усиливают шум. Поэтому применяют:
\begin{itemize}
    \item предварительное сглаживание;
    \item локальную полиномиальную аппроксимацию, например фильтр
          Савицкого--Голея;
    \item сглаживающие сплайны с последующим аналитическим дифференцированием;
    \item регуляризованные методы.
\end{itemize}
Порядок формулы сам по себе не гарантирует хорошего результата: широкие шаблоны
могут сильнее реагировать на шум и границы.

\subsection{Комплексный шаг}

Для функции, аналитической в окрестности вещественной оси и принимающей
вещественные значения при вещественном $x$ (то есть $\operatorname{Im}f(x)=0$),
при наличии комплексного продолжения
\[
    f'(x)=\frac{\operatorname{Im}f(x+\mathrm{i}h)}{h}+O(h^2).
\]
Метод не содержит вычитания близких чисел, поэтому допускает очень малые $h$.
Он неприменим, если вычислительный код не поддерживает комплексные аргументы или
содержит неаналитические операции.

\subsection{Что важно сказать на экзамене}

Нужно вывести формулы из рядов Тейлора, назвать их порядок и объяснить конфликт
между ошибкой усечения и округления. Численное дифференцирование существенно
чувствительнее к шуму, чем интегрирование, поэтому выбор шага и сглаживание имеют
принципиальное значение.

\newpage